\section{Ultraproducts}
\label{sec:types-ultraproducts-05}

Throughout this section, assume the Axiom of Choice in the ambient
metatheory.  It supplies an element of every Cartesian product of the
nonempty domains below and supplies the coordinatewise witness functions used
in the quantifier cases of \L o\'s's theorem.  This is a metatheoretic
assumption, not a first-order axiom of any factor structure.

Let \(\langle\mathcal M_i:i\in I\rangle\) be \(\Lang\)-structures and let
\(\mathcal U\) be an ultrafilter on \(I\).  On the Cartesian product
\(\prod_{i\in I}M_i\), define
\[
  a\sim_{\mathcal U}b
  \quad\Longleftrightarrow\quad
  \set{i\in I:a(i)=b(i)}\in\mathcal U.
\]

\begin{definition}[Ultraproduct]
The \emph{ultraproduct}
\[
  \prod_{i\in I}\mathcal M_i/\mathcal U
\]
has as its domain the set of \(\sim_{\mathcal U}\)-equivalence classes
\([a]_{\mathcal U}\).  Constants and functions are interpreted
coordinatewise.  For an \(n\)-ary relation symbol \(R\), define
\[
  R([a_1]_{\mathcal U},\ldots,[a_n]_{\mathcal U})
  \quad\Longleftrightarrow\quad
  \set{i\in I:\mathcal M_i\entails R(a_1(i),\ldots,a_n(i))}
  \in\mathcal U.
\]
When all \(\mathcal M_i\) are the same structure \(\mathcal M\), the result is
the \emph{ultrapower} \(\mathcal M^I/\mathcal U\).
\end{definition}

\begin{lemma}
The relation \(\sim_{\mathcal U}\) is an equivalence relation, and the
coordinatewise interpretations in the ultraproduct do not depend on the
chosen representatives.
\textit{Status: complete proof.}
\end{lemma}

\begin{proof}
Reflexivity and symmetry are immediate.  For transitivity, intersect the two
\(\mathcal U\)-large equality sets and use upward closure.  For
well-definedness, if representatives agree on \(\mathcal U\)-large sets, then
their common intersection is \(\mathcal U\)-large.  On that intersection the
corresponding function values agree, and the truth values of each relation
instance agree.  Upward closure completes both verifications.
\end{proof}

\begin{lemma}[Term evaluation in an ultraproduct]
\label{lem:los-terms}
Let \(t(x_1,\ldots,x_n)\) be a \(\Lang\)-term and let
\(a_1,\ldots,a_n\in\prod_{i\in I}M_i\).  Define
\(b_t\in\prod_{i\in I}M_i\) by
\[
  b_t(i)=t^{\mathcal M_i}(a_1(i),\ldots,a_n(i)).
\]
Then
\[
  t^{\prod_i\mathcal M_i/\mathcal U}
  ([a_1]_{\mathcal U},\ldots,[a_n]_{\mathcal U})
  =[b_t]_{\mathcal U}.
\]
\textit{Status: complete proof.}
\end{lemma}

\begin{proof}
Induct on \(t\).  For a variable \(x_j\), both sides are
\([a_j]_{\mathcal U}\).  For a constant, the claim is the coordinatewise
definition of its ultraproduct interpretation.  If
\(t=f(t_1,\ldots,t_m)\), apply the induction hypothesis to every \(t_j\) and
then use the coordinatewise definition of \(f\) in the ultraproduct.
\end{proof}

\begin{theorem}[\L o\'s's theorem]
\label{thm:los}
Assume the Axiom of Choice in the ambient set theory, as fixed for this
section.  Let \(\varphi(x_1,\ldots,x_n)\) be a first-order
\(\Lang\)-formula and let
\(a_1,\ldots,a_n\in\prod_{i\in I}M_i\).  Then
\[
  \prod_{i\in I}\mathcal M_i/\mathcal U
  \entails\varphi([a_1]_{\mathcal U},\ldots,[a_n]_{\mathcal U})
\]
iff
\[
  \set{i\in I:\mathcal M_i\entails
  \varphi(a_1(i),\ldots,a_n(i))}\in\mathcal U.
\]
\textit{Status: complete proof.}
\end{theorem}

The original source is \cite{Los1955Ultraproducts}.

\begin{proof}
We prove the displayed equivalence by structural induction for every formula
and every tuple of representatives for its free variables.  Brackets on a
tuple are coordinatewise, so
\([\bar a]_{\mathcal U}=([a_1]_{\mathcal U},\ldots,[a_n]_{\mathcal U})\).
For such a formula \(\theta\), write
\[
  X_\theta(\bar a)=
  \set{i\in I:\mathcal M_i\entails\theta(\bar a(i))}.
\]

If \(\theta\) is \(t=s\), \cref{lem:los-terms} and the definition of the
quotient give
\[
  \prod_i\mathcal M_i/\mathcal U\entails
  t([\bar a]_{\mathcal U})=s([\bar a]_{\mathcal U})
  \quad\Longleftrightarrow\quad
  \set{i:t^{\mathcal M_i}(\bar a(i))
               =s^{\mathcal M_i}(\bar a(i))}\in\mathcal U.
\]
For \(R(t_1,\ldots,t_m)\), the term lemma followed by the definition of the
ultraproduct relation gives the required equivalence.  For \(\bottom\), both
sides are false because \(X_{\bottom}=\varnothing\notin\mathcal U\).

For negation,
\(X_{\neg\alpha}=I\setminus X_\alpha\), and the ultrafilter alternative says
that this complement is large iff \(X_\alpha\) is not large.  For
conjunction,
\(X_{\alpha\land\beta}=X_\alpha\cap X_\beta\); a filter contains this
intersection iff it contains both factors.  For disjunction,
\(X_{\alpha\lor\beta}=X_\alpha\cup X_\beta\), and an ultrafilter contains a
finite union iff it contains at least one member.  Indeed, if it
contained neither, it would contain both complements and hence the complement
of the union.  Finally,
\[
  X_{\alpha\to\beta}=(I\setminus X_\alpha)\cup X_\beta.
\]
The union property and the ultrafilter alternative make this set large iff
\(X_\alpha\) is not large or \(X_\beta\) is large.  In each of these four
cases, the induction hypotheses and the corresponding truth clause in the
ultraproduct give the desired equivalence.

Suppose \(\theta(\bar x)=\exists y\,\psi(y,\bar x)\).  If the ultraproduct
has a witness \([b]_{\mathcal U}\), the induction hypothesis makes
\(X_\psi(b,\bar a)\) large.  This set is contained in
\(X_\theta(\bar a)\), so the latter is large.  Conversely, if
\(X_\theta(\bar a)\in\mathcal U\), choose for each index in that set a
witness \(b(i)\in M_i\), and choose any \(b(i)\in M_i\) at every remaining
index.  Then \(X_\psi(b,\bar a)\) contains \(X_\theta(\bar a)\) and is large;
the induction hypothesis makes \([b]_{\mathcal U}\) a witness in the
ultraproduct.

Suppose \(\theta(\bar x)=\forall y\,\psi(y,\bar x)\).  If
\(X_\theta(\bar a)\) is large, then for every \(b\in\prod_iM_i\) it is
contained in \(X_\psi(b,\bar a)\).  Upward closure and the induction
hypothesis show that every \([b]_{\mathcal U}\) satisfies \(\psi\), so the
ultraproduct satisfies \(\theta\).  Conversely, suppose the ultraproduct
satisfies \(\theta\) but \(X_\theta(\bar a)\notin\mathcal U\).  Its complement
is large.  At each index in that complement choose \(b(i)\) for which
\(\psi(b(i),\bar a(i))\) is false, and choose any value elsewhere.  If
\(X_\psi(b,\bar a)\) were large, its intersection with the large complement
of \(X_\theta(\bar a)\) would be empty and large, a contradiction.  Thus the
induction hypothesis says that \(\psi([b]_{\mathcal U},[\bar a]_{\mathcal U})\)
is false in the ultraproduct, contradicting its universal statement.  The two
quantifier cases use the stated Axiom of Choice and the stipulated
nonemptiness of every \(M_i\).
\end{proof}

\begin{corollary}[Diagonal elementary embedding]
\label{cor:diagonal-ultrapower}
The map \(d:M\to M^I/\mathcal U\) sending \(a\) to the class of the constant
function with value \(a\) is an elementary embedding.
\textit{Status: complete proof.}
\end{corollary}

\begin{proof}
For a formula with parameters from \(M\), its truth set under constant
representatives is either all of \(I\) or empty.  Apply
\cref{thm:los}.
\end{proof}

\begin{example}[An infinite ultraproduct of finite structures]
Let \(\mathcal M_n\) be the finite linear order
\(\{0,\ldots,n-1\}\), and let \(\mathcal U\) be a nonprincipal ultrafilter on
\(\N\setminus\{0\}\).  The ultraproduct is infinite: for each \(k\), the
sentence asserting that at least \(k\) distinct elements exist holds in all
\(\mathcal M_n\) with \(n\geq k\), a cofinite set of indices.  By
\cref{thm:los}, it holds in the ultraproduct.
\end{example}

\begin{warning}[First-order scope]
\L o\'s's theorem applies to first-order formulas.  The property ``the domain
is finite'' is not first-order definable over all structures, and the example
above must not be read as a contradiction.
\end{warning}
